\documentclass[book.tex]{subfiles}
\begin{document}
\chapter{Approximating the Schr\"{o}dinger Equation}

\label{chap:schrodinger_pinns}
\noindent\rule{11cm}{0.4pt} \\
\textbf{Prerequisites:}  \autoref{chap:pinns}, \autoref{chap:heat_pinns} \\
\textbf{Difficulty Level:} ** \\
\noindent\rule{11cm}{0.4pt}
\vspace{5mm}

In this chapter, we apply physics informed neural networks (PINNs) to model Schr\"{o}dinger's equation and test the ability of PINNs to handle periodic boundary conditions, complex-valued solutions, and different  nonlinearities in the governing partial differential equations. The nonlinear Schr$\mathrm{\ddot{o}}$dinger equation along with periodic boundary conditions is given by

\begin{align}
    i\hbar \frac{\partial}{\partial t} \psi &= \left[ -\frac{\hbar^2}{2m} \frac{\partial^2}{\partial x^2} + V(x) \right ] \psi
\end{align}
Let $h$ denote the solution wave function. Then we write:
\begin{align}
    i\hbar \frac{\partial}{\partial t} h &= -\frac{\hbar^2}{2m} \frac{\partial^2}{\partial x^2} h + V(x) h 
\end{align}

Moving terms over to the other side we get
\begin{align}
    i \hbar \frac{\partial}{\partial t} h + \frac{\hbar^2}{2m} \frac{\partial^2}{\partial x^2}  - V(x) h &= 0 \\
    i \frac{\partial}{\partial t} h + \frac{\hbar}{2m} \frac{\partial^2}{\partial x^2} h - \frac{V(x)}{\hbar} h &= 0 \\
\end{align}

To solve Schrodinger's equation with approximate methods, we need to specify a domain and initial conditions. Let $x \in [-5, 5]$, and $t \in [0,\pi/2]$. We impose initial conditions

\begin{align}
h(0,x) &= 2 \mathrm{sech}(x)\\
h(t,-5) &= h(t, 5)\\
\frac{\partial}{\partial x} h(t,-5) &= \frac{\partial}{\partial x} h(t, 5)
\end{align}

The final equation is 

\begin{align}
i \frac{\partial}{\partial t} h + \frac{1}{2} \frac{\partial^2}{\partial x^2} h + |h|^2 h &= 0\quad x \in [-5, 5],\quad t \in [0, \pi/2]
\end{align}

where $h(t,x)$ is the complex-valued solution. Let us define $f(t,x)$ to be given by

\begin{align}
f &= i \frac{\partial}{\partial t} h + \frac{1}{2} \frac{\partial^2}{\partial x^2} h + |h|^2 h,
\end{align}

and proceed by placing a complex-valued neural network prior on $h(t,x)$. In fact, if $u$ denotes the real part of $h$ and $v$ is the imaginary part, we are placing a multi-output neural network prior on 
\begin{align}
h(t,x) = \begin{bmatrix} u(t,x) & v(t,x)
\end{bmatrix}
\end{align}
This prior will result in the complex-valued (multi-output) physics informed neural network $f(t,x)$. The shared parameters of the neural networks $h(t,x)$ and $f(t,x)$ can be learned by minimizing the mean squared error loss

\begin{align}
    \textrm{MSE} = \textrm{MSE}_0 + \textrm{MSE}_b + \textrm{MSE}_f    
\end{align}
where
\begin{align}
\textrm{MSE}_0 &= \frac{1}{N_0}\sum_{i=1}^{N_0} |h(0,x_0^i) - h^i_0|^2\\
\textrm{MSE}_b &= \frac{1}{N_b}\sum_{i=1}^{N_b} \left(|h(t^i_b,-5) - h(t^i_b,5)|^2 + |h_x(t^i_b,-5) - h_x(t^i_b,5)|^2\right)\\ 
\textrm{MSE}_f &= \frac{1}{N_f}\sum_{i=1}^{N_f}|f(t_f^i,x_f^i)|^2.    
\end{align}

Here, $\{x_0^i, h^i_0\}_{i=1}^{N_0}$ denotes the initial data, and $\{t^i_b\}_{i=1}^{N_b}$ corresponds to the collocation points on the boundary, and $\{t_f^i,x_f^i\}_{i=1}^{N_f}$ represents the collocation points on $f(t,x)$. Consequently, $\textrm{MSE}_0$ corresponds to the loss on the initial data, $\textrm{MSE}_b$ enforces the periodic boundary conditions, and $\textrm{MSE}_f$ penalizes the Schr$\mathrm{\ddot{o}}$dinger equation not being satisfied on the collocation points.



% \end{example}
% The following figure summarizes the results of The experiment.


One potential limitation of the continuous time neural network models considered so far stems from the need to use a large number of collocation points $N_f$ in order to enforce physics informed constraints in the entire spatio-temporal domain. Although this poses no significant issues for problems in one or two spatial dimensions, it may introduce a severe bottleneck in higher dimensional problems, as the total number of collocation points needed to globally enforce a physics informed constrain (i.e., in this case a partial differential equation) will increase exponentially. In a later chapter, we introduce a different approach that circumvents the need for collocation points by introducing a more structured neural network representation leveraging classical Runge-Kutta time-stepping schemes.


%%%%%%%%%
%\section{Discrete-time PINNs for Schr\"{o}dinger's Equation}



%\begin{figure}
%    \centering
%    \includegraphics{figures/NLS_bookl.png}
%    \caption{Schrödinger equation: \textbf{Top:} Predicted
%solution along with the initial and boundary training data. In addition
%we are using 20,000 collocation points generated using a Latin Hypercube
%Sampling strategy. \textbf{Bottom:} Comparison of the predicted and
%exact solutions corresponding to the three temporal snapshots depicted
%by the dashed vertical lines in the top panel.}
%\end{figure}


\end{document}